\section{A quadratic Chabauty entry through the first quotient}
\label{qh-section}

We give the now independently reviewed ordinary application of the
quadratic Chabauty finiteness criterion. No Chabauty set is computed in
this section. Put \(D=D_{3,\zeta}\), let \(H_j\) be the three curves of
Section~\ref{qs-section}, and let \(J_D=\operatorname{Jac}(D)\).
The established decompositions and exact elliptic rank give
\[
 J_D\sim_{\mathbb Q} E^2\times J_2\times J_3,
 \qquad J_1\sim_{\mathbb Q}E^2,
 \qquad E:Y^2=X^3-9X-9,\quad\operatorname{rank}E(\mathbb Q)=1.
\]
The full projective morphism \(\pi:D\to H_1\) has degree two. On its
dense chart it is \((s,v,w)\mapsto(s,y=Bv)\), where \(B=f_1(s)\).

\begin{lemma}[Rational divisor classes and good reduction]\label{qh-hypotheses}
For the Neron--Severi rank over the ground field,
\[
 \rho_{\mathbb Q}(J_1)\ge3,\qquad \rho_{\mathbb Q}(J_D)\ge5.
\]
Both \(H_1\) and \(D\) have rational basepoints and good reduction at five.
\end{lemma}
\begin{proof}
On \(E^2\), the rational divisor classes of \(E\times\{O\}\),
\(\{O\}\times E\) and the diagonal have intersection matrix
\[
 \begin{pmatrix}0&1&1\\1&0&1\\1&1&0\end{pmatrix},
 \qquad\det=2.
\]
They are independent modulo algebraic equivalence since algebraically
trivial divisors are numerically trivial. Hence \(\rho_{\mathbb Q}(E^2)\ge3\).
Pullback by an isogeny is injective on Neron--Severi groups tensored with
\(\mathbb Q\): the norm of a pulled-back line bundle is its
degree-th tensor power, and norm preserves algebraic equivalence.
There are isogenies in both directions, so the rational ranks agree.
This proves the first inequality. On \(E^2\times J_2\times J_3\),
add a rational ample class from each \(J_j\). Such classes exist by
projectivity over \(\mathbb Q\). Restriction to the three factors proves
the independence of all five classes, giving the second inequality.
No assertion of equality, absolute Picard rank, or absence of CM is needed.

The smoothness argument of Theorem~\ref{qs-five} applies to \(f_0f_1\)
as well: both discriminants are units at five, their difference is
\(3s(s+1)\), and neither vanishes at \(0,-1\). The two infinity points
with leading ordinate \(\pm1\) are smooth; either is a rational basepoint.

For \(D\), the three branch divisors on \(\mathbb P^1\) are the zeros
of \(f_0,f_1,f_4\). Each is finite etale over \(\mathbb Z_5\) and
they are disjoint. Normalize in the biquadratic extension
\[
 v^2=f_0/f_1,\qquad w^2=f_4/f_1.
\]
Away from these divisors the cover is etale. At a zero of \(f_0\) or
\(f_4\), only one square root ramifies. Etale-locally on the base its
normalized equation is \(z^2=t\) times a unit; it is smooth over
\(\mathbb Z_5\) because the coefficient of the local parameter \(t\)
is a unit. The other square root is an etale unit extension. At a zero
of \(f_1\), use \(z=1/v\) and \(w/v\): the first gives the same simple
ramification and the latter has square a unit. Thus the normalization
is smooth at every branch. It is proper and finite over
\(\mathbb P^1_{\mathbb Z_5}\) and has the smooth genus-six generic fibre,
so it supplies good reduction. Infinity is unramified. The point
\((s,v,w)=(0,1,1)\) is rational and nonbranch and supplies a basepoint.
\end{proof}

\begin{theorem}[A finite five-adic container through the first quotient]
\label{qh-container}
Let \(H_1(\mathbb Q_5)_2\) denote the quadratic Chabauty set of
Balakrishnan--Dogra. Then
\begin{gather*}
 H_1(\mathbb Q_5)_2\text{ is finite},\\
 \mathcal C_5=\{P\in D(\mathbb Q_5):\pi(P)\in H_1(\mathbb Q_5)_2\}
 \text{ is finite},\\
 D(\mathbb Q)\subseteq\mathcal C_5.
\end{gather*}
\end{theorem}
\begin{proof}
The exact GD rank is \(\operatorname{rank}J_1(\mathbb Q)=2\), while
the preceding lemma gives
\[
 2<2+\rho_{\mathbb Q}(J_1)-1,
\]
whose right side is at least four. All hypotheses of
\cite[Lemma~3.2]{BalakrishnanDograQCI} apply: the field is \(\mathbb Q\),
the curve is smooth geometrically integral and projective of genus
greater than one, it has a rational basepoint, and five is a good
reduction prime and splits completely in the base field. The Picard
number in that criterion is precisely over \(\mathbb Q\).
It follows that \(H_1(\mathbb Q_5)_2\) is finite; its defining Selmer
condition includes \(H_1(\mathbb Q)\).

The projective quotient \(\pi\) is finite of degree two. Every fibre
over a \(\mathbb Q_5\)-point has at most two geometric points, including
ramified fibres and all chart exceptions. Its inverse image of the
finite set is therefore finite. Every rational point maps to a rational
point of \(H_1\), proving the asserted inclusion.
\end{proof}

This theorem does not compute \(\mathcal C_5\), its cardinality or
coordinates, or determine which of its points are rational. It does not
identify \(\mathcal C_5\) with the intrinsic set \(D(\mathbb Q_5)_2\).
The quotient entry avoids an upper rank bound for \(J_2,J_3\) for this
container. It does not classify the positive common-source locus.
An actual calculation still needs rigorous heights and integrals,
finite local-height value sets, sufficient precision, and a rational-point
or Mordell--Weil sieve. A finite search for multiples of a displayed
elliptic point supplies none of these completeness assertions; GD by
itself does not prove that point to be a global generator.

\begin{proposition}[A separate conditional entry on the full curve]
\label{qh-direct-conditional}
If the additional bound
\[
 \operatorname{rank}J_2(\mathbb Q)+\operatorname{rank}J_3(\mathbb Q)\le6
\]
holds, then \(D(\mathbb Q_5)_2\) itself is finite.
\end{proposition}
\begin{proof}
The bound would give \(\operatorname{rank}J_D(\mathbb Q)\le8\), while
\[
 \operatorname{genus}(D)+\rho_{\mathbb Q}(J_D)-1\ge6+5-1=10.
\]
The other source hypotheses were verified in Lemma~\ref{qh-hypotheses};
the same finiteness criterion applies.
\end{proof}
The extra upper bound is unproved. The positive even lower bounds from
Theorem~\ref{qs-rank-lower} mean that it would allow total genus-six
Jacobian rank six or eight. If the total rank were ten or greater,
the particular estimate \(\rho\ge5\) would simply not establish this
criterion; larger Picard number or other constructions remain possible.

\paragraph{Source and scope.}
The setup and the exact criterion were checked in
\cite[printed p.~3, Section~2, and Lemmas~3.1--3.2 on printed p.~12]{BalakrishnanDograQCI}.
They concern general smooth projective curves, not only hyperelliptic
ones, and use the ground-field Picard number. The explicit height formula
in that paper's Theorem~1.2 has additional assumptions, including
rank equal to genus and a finite-index closure of the rational Jacobian
group in the local group. Those algorithmic assumptions are not claimed
or needed by the abstract finiteness argument of this section. The
intersection calculation and divisor norm argument above prove the
stated Neron--Severi bounds. No Chabauty, Neron--Severi, rank or complete
point-set theorem is claimed here as a Lean formalization.
